% \documentclass[11pt]{article}

% %%%% CUSTOM PREAMBLE
% \usepackage{./preamble}

% %% ADD MISC DEFINITIONS HERE %%%%%%%%%%%%%%%%%%%%%%%%%%%%%%%%%%%

% %%%%%%%%%%%%%%%%%%%%%%%%%%%%%%%%%%%%%%%%%%%%%%%%%%%%%
% %\addtolength{\topmargin}{-1cm}
% %\addtolength{\textheight}{2cm}

% \begin{document}

%\maketitle

\newpage
 %%*****************************************************
\noindent
\textbf{CS4268 AY2025-26 Sem 2}
\\
\noindent
Lecturer: Divesh Aggarwal
\hfill
Lecture 6                      %%% FILL IN LECTURE NUMBER HERE
\hfill
Date: February 19, 2026          %%% FILL IN LECTURE DATE HERE


\noindent
\rule{\textwidth}{1pt}

\medskip

%%%%%%%%%%%%%%%%%%%%%%%%%%%%%%%%%%%%%%%%%%%%%%%%%%%%%%%%%%%%%%%%
%% BODY OF SCRIBE NOTES GOES HERE
%%%%%%%%%%%%%%%%%%%%%%%%%%%%%%%%%%%%%%%%%%%%%%%%%%%%%%%%%%%%%%%%

%%%Write lecture topic above.
\textbf{Instructions}\\
\begin{itemize}
    \item Try to leave the preamble unchanged, only make changes to \texttt{main.tex} and  \texttt{references.bib}. Add images to the \texttt{images} folder. 
    \item Share one (zipped) directory preserving the directory structure on Canvas. 
    \item Prepare organized notes by elaborating on the broad items provided in different sections below based on the lecture material completing any missing details from lectures.
    \item Feel free to modify structure and/or order of contents as you see fit; the below is just a guideline.
\end{itemize} 


\section{Foundational Quantum Algorithms}
\kk{For all material here, add/remove details as you see fit. Please add quantum circuits if available, use the quantikz package.}

\subsection{Quantum Fourier Transform}
\begin{enumerate}
    \item Binary representation: if $\alpha \geq 1$, $\alpha = \alpha_1\dots\alpha_n = \alpha_12^{n-1}+\dots+\alpha_n2^0$; if $0 \leq \alpha < 1$, $\alpha = 0.\alpha_1\dots\alpha_n = \alpha_12^{-1}+\dots+\alpha_n2^{-n}$. Here, $\alpha_i \in\{0,1\}$ for all $i$. Note that if $\alpha = \alpha_1\dots\alpha_n$, then $\alpha/2^n = 0.\alpha_1\dots\alpha_n$.
    
    \item Let $\mathbf{x}\in \mathbb{C}^N$ (recall for simplicity we take $N=2^n$ for some $n \in \mathbb{N}$). The DFT (discrete Fourier transform) implements $\mathbf{x} \mapsto \mathbf{y}=\text{DFT}(\mathbf{x})$, where
    \begin{align*}
        DFT = \frac{1}{\sqrt{N}} \sum_{j,k=0}^{N-1} e^{2\pi ijk/N} \ketbra{j}{k}.
    \end{align*}
    The QFT implements the same transformation: for a basis state $\ket{j}$ we have
    \begin{align*}
        \mathsf{QFT}\ket{j} = \frac{1}{\sqrt{N}} \sum_{k=0}^{N-1} e^{2\pi ijk/N} \ket{k}.
    \end{align*}
    In binary, this is
    \begin{align*}
        \mathsf{QFT}\ket{j_1\dots j_n} = \frac{1}{\sqrt{2^n}}\left( \ket{0}+e^{2\pi i0.j_n}\ket{1} \right)\left( \ket{0}+e^{2\pi i0.j_{n-1}j_n}\ket{1} \right)\cdots \left( \ket{0}+e^{2\pi i0.j_1j_2\dots j_n}\ket{1}\right).
    \end{align*}
    
    \item Fast Fourier Transform: Cooley-Tukey Algorithm. Start from 
    \begin{align*}
        y_k &= \sum_{j=0}^{N-1} e^{\frac{2\pi i jk}{N}} x_j\\
        &= \sum_{j=0}^{N/2-1} e^{\frac{2\pi i (2j)k}{N}} x_{2j} + \sum_{j=0}^{N/2-1} e^{\frac{2\pi i (2j+1)k}{N}} x_{2j+1}\\
        &= \underbrace{\sum_{j=0}^{N/2-1} e^{\frac{2\pi i jk}{N/2}} x_{2j}}_{E_k} + e^{\frac{2\pi ik}{N}}\underbrace{\sum_{j=0}^{N/2-1} e^{\frac{2\pi i jk}{N/2}} x_{2j+1}}_{O_k}.
    \end{align*}
    For brevity we have absorbed $1/\sqrt{N}$ into the $x_i$. The key is to notice that 
    \begin{align*}
        y_{k+\frac{N}{2}} = E_k - e^{\frac{2\pi ik}{N}}O_k
    \end{align*}
    so we only have to compute $y_k$ from $k=0$ until $k=N/2-1$. We recursively break this up into smaller DFTs until we reach the base case of $N=2$. Essentially, the key breakthrough of Cooley-Tukey is to reuse results of intermediate computations to compute the $y_k$'s. One can see total complexity is $O(N \log N)$ by drawing the binary tree expansion of this procedure: at each level $l$, there are $2^l$ $E_k/O_k$ terms, and for each such term we have $N/2^l$ $k$'s. Each $E_k/O_k$ term is the sum of its children nodes, so takes $O(1)$ computations to evaluate. So each level takes $O(N)$ operations, and there are $\log N$ levels.

    \item Derivation of binary representation of QFT:
    \begin{align*}
        \mathsf{QFT}\ket{j_1\dots j_n} &= \frac{1}{\sqrt{2^n}} \sum_{k=0}^{2^n-1} e^{\frac{2\pi ijk}{2^n}} \ket{k}\\
        &= \frac{1}{\sqrt{2^n}} \sum_{k_1,\dots k_n=0}^{1} e^{2\pi ij(\sum_{l=1}^n k_l 2^{-l})} \ket{k_1 \dots k_n}\\
        &= \frac{1}{\sqrt{2^n}} \sum_{k_1,\dots k_n=0}^{1} \bigotimes_{l=1}^n e^{2\pi ijk_l 2^{-l}} \ket{k_l}\\
        &= \frac{1}{\sqrt{2^n}} \bigotimes_{l=1}^n \left[\sum_{k_l=0}^{1} e^{2\pi ijk_l 2^{-l}} \ket{k_l}\right]\\
        &= \frac{1}{\sqrt{2^n}} \bigotimes_{l=1}^n \left[\ket{0} + e^{2\pi ij 2^{-l}} \ket{1}\right]\\
        &= \frac{1}{\sqrt{2^n}}\left( \ket{0}+e^{2\pi i0.j_n}\ket{1} \right)\left( \ket{0}+e^{2\pi i0.j_{n-1}j_n}\ket{1} \right)\cdots \left( \ket{0}+e^{2\pi i0.j_1j_2\dots j_n}\ket{1}\right).
    \end{align*}
    This representation gives intuition toward constructing an efficient circuit for the QFT. With the phase gates $R_k := \text{diag}(1,e^{\frac{2\pi i}{2^k}})$, and $M_i$ meaning single-qubit gate $M$ acts on qubit $i$, the circuit comprises first Hadamarding then applying ctrl-$R_k$'s, see \cite{NC10} pg219 or wiki.

    \item Speedup: FFT takes $\Theta(N \log N)/\Theta(n2^n)$, QFT takes $\Theta(\log^2 N)/\Theta(n^2)$. Caveat: the Fourier-transformed output $\mathsf{QFT}\ket{x}$ is not $\text{DFT}(\mathbf{x})$ itself, rather $\mathsf{QFT}\ket{x}$ has the entries of $\text{DFT}(\mathbf{x})$ encoded in its amplitudes. To extract the values of the amplitudes we generally have to perform measurements which generally consumes more resources.
\end{enumerate}


\subsection{The Deutsch-Jozsa Family}

First consider the baby version, `Deutsch'. `Deutsch-Jozsa' is a generalization of Deutsch to $n$-input bit functions.

Deutsch
\begin{enumerate}
    \item Idea: Determine whether a one-bit function $f$ is constant or balanced. $f$ is implemented through the oracle $O_f\ket{x,y} = \ket{x,y \oplus f(x)}$ which is a permutation and hence also a unitary.
    \item Method: Start from $\ket{01}$, apply $O_f \circ H \otimes H$ to get (up to a global phase) $\ket{\pm}\ket{-}$ ($+$ if $f(0)=f(1)$, $-$ if $f(0) \neq f(1)$). Discard second qubit, Hadamard the first, so we get $\ket{0} \iff f(0)=f(1) \iff f(0)\oplus f(1) = 0$ (constant), $\ket{1} \iff f(0)\neq f(1) \iff f(0)\oplus f(1) = 1$ (balanced). Only requires one $O_f$ call.
    \item Note that key is $O_f\ket{x}\ket{-} = (-1)^{f(x)}\ket{x}\ket{-}$, `phase-kickback'.
\end{enumerate}

Deutsch-Jozsa: Generalization of Deutsch to $n$-bit functions. Leverages massively parallel phase-kickback.
\begin{enumerate}
    \item Note that $H^{\otimes n}\ket{x}=\frac{1}{\sqrt{2^n}}\sum_{y}(-1)^{x\cdot y}\ket{y}$.
    \item Start from $\ket{0^n}\ket{1}$, apply $H^{\otimes n} \circ O_f \circ H^{\otimes n}\otimes H$ to get the state
    \begin{align*}
        \sum_{z \in \{0,1\}^n} \left( \frac{1}{2^n} \sum_{x \in \{0,1\}^n} (-1)^{f(x)+x \cdot z} \right) \ket{z}.
    \end{align*}
    Then measure the working $n$ qubits. Observe that $P(z=0^n) = |\frac{1}{2^n} \sum_{x \in \{0,1\}^n} (-1)^{f(x)}|^2$ equals 1 iff $f$ is constant and equals 0 iff $f$ is balanced. Again, only requires one $O_f$ call.
    \item Classically, deterministic worst case needs $2^{n-1}+1=O(2^n)$ queries. Randomized algorithm: randomly pick $k$ distinct inputs $x_i$ and query to get $f(x_i)$, answer constant if $f(x_1)=\dots=f(x_k)$, balanced otherwise. If $f$ was really constant then method works with certainty, if $f$ was balanced then $f(x_1)=\dots=f(x_k)$ occurs with probability $\frac{1}{2^{k-1}}$. So if probability-error tolerance is $\varepsilon$, need $k=O(\log \frac{1}{\varepsilon})$.
\end{enumerate}

Related to DJ is Bernstein-Vazirani, albeit with a slightly different end goal. Algorithm similar.

Bernstein-Vazirani: Given $f: \{0,1\}^n \rightarrow \{0,1\}$ s.t. $f(x)=x \cdot s$ for some unknown $s$, find $s$.
\begin{enumerate}
    \item Classically, most efficient method is to obtain $f(1,0,\dots)=s_1, f(0,1,\dots)=s_2, \dots$. This requires $n$ queries.
    \item Quantumly, the algorithm is the same as DJ: apply $H^{\otimes n} \circ O_f \circ H^{\otimes n}\otimes H$. Miraculously, the final output is $\ket{s}$!
    \item Ok actually nothing miraculous - it's simply because after applying the query, we get the state $\frac{1}{\sqrt{2^n}} \sum_x (-1)^{x \cdot s} \ket{x}$, which is $H^{\otimes n}\ket{s}$. And Hadamards are self-inverse.
\end{enumerate}

\subsection{Simon's Problem and Algorithm}

\begin{enumerate}
    \item Problem: given $f: \{0,1\}^n \rightarrow \{0,1\}^n$ s.t. $f(x)=f(y) \iff x=y\oplus s$, find $s$. 
    \item Classical Case
    \begin{enumerate}
        \item A randomized algorithm: randomly pick $k$ (TBD) distinct inputs $x_i$ and query to get $f(x_i)$. If there is no collision, answer $s=0^n$. If there is a collision $f(x_i)=f(x_j)$, then answer $s=x_i \oplus x_j$. 
        \item Analysis: If actually $s=0^n$ then the algorithm works with certainty. If actually $s \neq 0^n$ then we might answer wrongly if there is no observed collisions. Let the probability that the samples have no collision after $k\geq 2$ queries be $P(k)$. Suppose after $j-1$ queries there is no matching. The probability that on the $j$th, i.e. next query there is still no collision with one of the previous queries is then $1-\frac{j-1}{2^n-(j-1)}$. Thus after $k$ rounds,
        \begin{align*}
            P(k) &= \prod_{j=2}^k \left[ 1- \frac{j-1}{2^n-(j-1)} \right]\\
            &\geq 1-\sum_{j=2}^k \left[ \frac{j-1}{2^n-(k-1)} \right]\\
            &\geq 1- \frac{\Theta(k^2)}{2^n-\Theta(k)}.
        \end{align*}
        Remember that we \textit{do} want a collision! So we want $P(k)$ to be small, say $\varepsilon$ (equivalently, coll. prob. is $1-\varepsilon$). Then $P(k)<\varepsilon \implies k \in \Omega(2^{n/2})$. Here in the first $\geq$ we used $(1-a)(1-b)\geq 1-(a+b)$ for $0\leq a,b \leq 1$. Also note that we can assume from the outset that $k\leq 2^{n/2}$ because if $k>2^{n/2}$ then there has to be a collision.
        \item Simon also showed \textit{any} randomized algorithm needs $k=\Omega(2^{n/2})$ queries (proof is similar to the one above).
    \end{enumerate}
    
    \item Simon's Quantum Algorithm
    \begin{enumerate}
        \item Quantum part: Start with $n$ work qubits and $n$ ancilla qubits $\ket{0^n}\ket{0^n}$. Apply $(H^{\otimes n}\otimes I)O_f(H^{\otimes n} \otimes I)$ then measure the work qubits.
        \item Classical postprocessing: repeat $n-1$ times to get $n-1$ strings $y_i$ all satisfying $y_i \cdot s =0$, where $y_i$ are linearly independent with positive probability. If lin. ind., use Gaussian elimination to solve system, giving two candidates $0^n$ and $s' \neq 0^n$. 
        \item Test: if $f(0^n)=f(s')$, answer $s=s'$; otherwise $f(0^n)\neq f(s')$, answer $s=0^n$.
        \item Boosting: repeat entire process $k$ times to boost prob. of obtaining lin. ind. $y_i$. (LI can be easily checked.) If for \text{any} run $f(0^n)=f(s')$, answer $s=s'$; otherwise if for \text{all} runs $f(0^n)\neq f(s')$, answer $s=0^n$.
    \end{enumerate}
    
    \item Analysis of Simon's algorithm
    \begin{enumerate}
        \item Quantum: after applying the unitaries the state is
        \begin{align*}
            \ket{\psi} = \sum_y \ket{y} \left( \frac{1}{2^n} \sum_x (-1)^{x \cdot y} \ket{f(x)} \right).
        \end{align*}
        If $s=0^n$, $f$ is injective so $p(y) = \frac{1}{2^n}$ for all $y$. If $s \neq 0^n$, $f$ is two-to-one, so if $z \in \text{range}(f)$ then $f(x_z) = f(x_z') = z$ for $x_z = x_z' \oplus s$. Then $p(y) = \frac{1}{2^{n-1}}$ if $y \cdot s = 0$, and $p(y) = 0$ if $y \cdot s = 1$.
        \item Postprocessing: suppose we have lin. ind. $\{y_i\}_{i=1}^{n-1}$. GE gives the two possible solutions: $0^n,s'\neq 0^n$. Suppose actually $s=0^n$. Then $f(0^n)\neq f(s')$, so our answer $s=0^n$ is always correct. Suppose actually $s \neq 0^n$. Then $f(s)=f(0^n)=f(s')$, so our answer $s'=s$ is always correct.
        \item Probability of linear independence: given $m$ LI strings $y_1,\dots,y_m$, spanning subspace is size $2^m$. Prob of $y_{m+1}$ being LI to this is $\frac{2^n-2^m}{2^n}$. So prob of $n-1$ LI strings is $P=(1-\frac{2^1}{2^n})\dots(1-\frac{2^{n-2}}{2^n}) \geq 1-\sum_{i=1}^{n-2} \frac{1}{2^i} > \frac{1}{2}$. So prob. of \textit{no LI} even after running $k$ times is at most $(1-\frac{1}{2})^k < e^{-k/2}$. So prob. of \textit{at least one LI} after running $k$ times is at least $1 - e^{-k/2}$.
    \end{enumerate}
\end{enumerate}



















% \bibliographystyle{alpha}
% \bibliography{main.bib}
% \end{document}

